\subsection{Boundary weights on a quotient surface}
\label{subsec:quotient-boundary-ideals}

The resolution-ideal relation above uses smooth volume forms on smooth models. An integrability ideal on a singular space can instead use a volume form determined by its canonical bundle and a boundary divisor. The following local example shows how this additional weight changes the ideal. In this subsection, replace the compact space fixed above by the affine surface
\[
 X=\{z_1z_2=z_3^2\}\subset\mathbb C^3,
 \qquad E=\{z_1=z_3=0\},\qquad D=\tfrac12 E.
\]
The finite map
\[
 \pi\colon\mathbb C^2\longrightarrow X,
 \qquad \pi(s,t)=(s^2,t^2,st),
\]
identifies $X$ with the quotient by $(s,t)\mapsto(-s,-t)$. It has degree two and is unramified away from the origin. Holomorphic germs at the origin of $X$ correspond to convergent series in $s,t$ containing only monomials $s^at^b$ with $a+b$ even. Indeed, every such series has the form $F_0(s^2,t^2)+stF_1(s^2,t^2)$ with $F_0,F_1$ convergent, which is the pullback of $F_0(z_1,z_2)+z_3F_1(z_1,z_2)$.

On $X_{\mathrm{reg}}$, the canonical two-form
\[
 \Omega=\frac{\mathrm{d}z_1\wedge\mathrm{d}z_3}{z_1}
       =\frac{\mathrm{d}z_3\wedge\mathrm{d}z_2}{z_2}
\]
is defined on the respective charts $z_1\neq0$ and $z_2\neq0$. These charts cover $X_{\mathrm{reg}}$, the formulas agree, and $\pi^*\Omega=2\,\mathrm{d}s\wedge\mathrm{d}t$. The residue form $\Omega$ generates the dualizing sheaf of this hypersurface. Since $\operatorname{div}(z_1)=2E$, the meromorphic fourth tensor power
\[
 \beta=\frac{\Omega^{\otimes4}}{z_1}
\]
generates $\omega_X^{[4]}(4D)$, the reflexive fourth power of the dualizing sheaf twisted by $4D$. The fourth root of its squared absolute value defines the \emph{adapted volume} $\mathrm{d}V_D=|\beta|^{2/4}=|z_1|^{-1/2}|\Omega|^2$ on $X_{\mathrm{reg}}$; for zero boundary we use $\mathrm{d}V_0=|\Omega|^2$.\index{Adapted volume!quotient surface} Their pullbacks are
\begin{equation}\label{eq:quotient-adapted-volumes}
 \pi^*\mathrm{d}V_D=C_D|s|^{-1}\,\mathrm{d}A_s\,\mathrm{d}A_t,
 \qquad
 \pi^*\mathrm{d}V_0=C_0\,\mathrm{d}A_s\,\mathrm{d}A_t,
\end{equation}
where $\mathrm{d}A$ denotes Euclidean area and $C_D,C_0>0$ depend only on the convention for absolute values of forms. Replacing a generator by a holomorphic unit multiplies the volume by a positive factor bounded above and below on smaller neighborhoods.

For a boundary $D'\in\{0,D\}$ and a plurisubharmonic function $u$ on this $X$, define the adapted-volume integrability ideal by
\begin{equation}\label{eq:quotient-adapted-ideal}
 \mathcal I_{D'}^{\mathrm{ad}}(u)_x
 =\left\{f\in\mathcal O_{X,x}:
   \int_{U\cap X_{\mathrm{reg}}}|f|^2\mathrm{e}^{-u}\,\mathrm{d}V_{D'}<\infty\right\},
\end{equation}
where the neighborhood $U$ of $x$ may be chosen for each germ $f$. This uses the weight convention of \cref{def:multiplier-ideal}. It defines a different object from the resolution-ideal class: the volume in \eqref{eq:quotient-adapted-ideal}, including its boundary weight, is part of the definition.

In logarithmic polar coordinates, integrability reduces to the following double integral.

\Needspace{12\baselineskip}
\begin{lemma}\label{lem:weighted-logarithmic-maximum-integral}
Let $A,B,\alpha,\beta,T>0$. Then
\[
 \int_T^\infty\int_T^\infty
  \exp\bigl(-\alpha p-\beta q+\min(Ap,Bq)\bigr)
       \,\mathrm{d}q\,\mathrm{d}p<\infty
 \quad\Longleftrightarrow\quad
 \frac\alpha A+\frac\beta B>1.
\]
\end{lemma}
\begin{proof}
Split the quadrant along $q=Ap/B$. On the part $q\geq Ap/B$, set $p_0=\max\{T,BT/A\}$. After discarding the strip $T<p<p_0$, integration in $q$ gives
\[
 \frac1\beta
 \int_{p_0}^\infty
   \exp\left(-\left[\alpha-A+\frac{\beta A}{B}\right]p\right)
      \,\mathrm{d}p.
\]
The discarded strip has finite integral because $\beta>0$. On the other part, integration in $p$ gives the corresponding tail with coefficient $\beta-B+\alpha B/A$; its discarded strip is finite because $\alpha>0$. The two coefficients are $A$ and $B$ times $\alpha/A+\beta/B-1$, and a tail with a nonpositive coefficient diverges.
\end{proof}

\Needspace{15\baselineskip}
\begin{proposition}\label{prop:quotient-boundary-ideals}\index{Multiplier ideal!adapted volume}\index{Unibranch space!boundary weights}\index{Quotient surface!multiplier ideals}
The plurisubharmonic function
\[
 u=\max\{\sqrt2\log|z_1|^2,\,2\log|z_2|^2\}
\]
on $X$ does not have analytic singularities at the origin. Its adapted-volume ideals are
\[
 \mathcal I_D^{\mathrm{ad}}(u)=(z_1,z_3,z_2^2)\mathcal O_X,
 \qquad
 \mathcal I_0^{\mathrm{ad}}(u)=(z_1,z_2,z_3)\mathcal O_X.
\]
In particular, the same potential requires the generator $z_2^2$ in place of $z_2$ when the boundary $D$ is included.
\end{proposition}
\begin{proof}
The function is the restriction of a maximum of two plurisubharmonic functions on $\mathbb C^3$. On the axes parametrized by $(z_1,z_2,z_3)=(\zeta,0,0)$ and $(0,\zeta,0)$, its logarithmic coefficients are $\sqrt2$ and $2$. An analytic singularity $c\log\sum_j|f_j|^2+O(1)$ restricts to a logarithmic coefficient equal to $c$ times an integer on either axis, provided that it is finite off the origin there. The ratio of the two coefficients would therefore be rational, which it is not.

Pullback by the finite degree-two map reduces local integrability to the cover. Neighborhoods of the origin contain and are contained in images of sufficiently small bidisks, so it suffices to integrate over a bidisk in $s,t$. The pullback weight is
\[
 u\circ\pi=\max\{2\sqrt2\log|s|^2,\,4\log|t|^2\}.
\]
For $s^at^b$, $a,b\geq0$, set $p=-\log|s|^2$ and $q=-\log|t|^2$. By \eqref{eq:quotient-adapted-volumes}, angular integration over the bidisk of radius $\mathrm e^{-T/2}$ gives, up to a positive constant,
\[
 \int_T^\infty\int_T^\infty
  \exp\left(-\left[a+\tfrac12\right]p-(b+1)q
             +\min(2\sqrt2\,p,4q)\right)
       \,\mathrm{d}q\,\mathrm{d}p
\]
for the boundary $D$. For zero boundary, replace $a+\tfrac12$ by $a+1$. Thus \cref{lem:weighted-logarithmic-maximum-integral} gives the respective conditions
\begin{equation}\label{eq:quotient-monomial-tests}
 \frac{a+\tfrac12}{2\sqrt2}+\frac{b+1}{4}>1,
 \qquad
 \frac{a+1}{2\sqrt2}+\frac{b+1}{4}>1.
\end{equation}

These monomial tests determine the entire ideal. Indeed, the measure and the weight on the cover are invariant under separate rotations of $s,t$. Parseval's identity and Tonelli's theorem show that every nonzero coefficient of a weighted-integrable holomorphic series must satisfy its monomial test. Only exponents with $a+b$ even occur.

\Needspace{6\baselineskip}
In the first test, the minimal allowed exponents, under addition of nonnegative even-total exponents, are
\[
 (2,0),\quad(1,1),\quad(0,4).
\]
To see this, every term with $a\geq2$ is divisible by $s^2$, every term with $a=1$ has odd $b\geq1$ and is divisible by $st$, and a term with $a=0$ has even $b$ and passes the test exactly when $b\geq4$. Each of the three displayed exponents passes the test. Splitting a convergent invariant series into these three groups and dividing gives convergent invariant coefficient series on a smaller bidisk. They descend to holomorphic germs on $X$, bounded on a smaller neighborhood. Since the three generators are weighted-integrable, every such combination is weighted-integrable. This proves the first ideal at the origin. For the second test, every nonconstant even-total monomial is allowed, with minimal exponents $(2,0),(1,1),(0,2)$, and the same argument gives the second ideal.

Away from the origin, $u$ is locally bounded and $X$ is smooth. The only additional volume singularity lies along $E$: with $z_2\neq0$, $z_3$ is a normal coordinate and $z_1=z_3^2/z_2$, so $\mathrm{d}V_D$ is a smooth positive factor times $|z_3|^{-1}$ times smooth volume. This is locally integrable. Both ideals therefore equal $\mathcal O_X$ there, as do the ideals generated by the displayed coordinates, which proves the sheaf identities.
\end{proof}

The calculation depends on the prescribed volume as well as on $u$. It does not change the resolution-ideal equivalence used in \cref{prop:space-resolution-i-envelope}. The next calculation expresses the exponent test for a general full-dimensional affine toric cone.
